% Options for packages loaded elsewhere
\PassOptionsToPackage{unicode}{hyperref}
\PassOptionsToPackage{hyphens}{url}
%
\documentclass[
  11pt,
  a4paper,
]{article}
\usepackage{amsmath,amssymb}
\usepackage{iftex}
\ifPDFTeX
  \usepackage[T1]{fontenc}
  \usepackage[utf8]{inputenc}
  \usepackage{textcomp} % provide euro and other symbols
\else % if luatex or xetex
  \usepackage{unicode-math} % this also loads fontspec
  \defaultfontfeatures{Scale=MatchLowercase}
  \defaultfontfeatures[\rmfamily]{Ligatures=TeX,Scale=1}
\fi
\usepackage{lmodern}
\ifPDFTeX\else
  % xetex/luatex font selection
\fi
% Use upquote if available, for straight quotes in verbatim environments
\IfFileExists{upquote.sty}{\usepackage{upquote}}{}
\IfFileExists{microtype.sty}{% use microtype if available
  \usepackage[]{microtype}
  \UseMicrotypeSet[protrusion]{basicmath} % disable protrusion for tt fonts
}{}
\makeatletter
\@ifundefined{KOMAClassName}{% if non-KOMA class
  \IfFileExists{parskip.sty}{%
    \usepackage{parskip}
  }{% else
    \setlength{\parindent}{0pt}
    \setlength{\parskip}{6pt plus 2pt minus 1pt}}
}{% if KOMA class
  \KOMAoptions{parskip=half}}
\makeatother
\usepackage{xcolor}
\usepackage[margin=2.4cm]{geometry}
\usepackage{color}
\usepackage{fancyvrb}
\newcommand{\VerbBar}{|}
\newcommand{\VERB}{\Verb[commandchars=\\\{\}]}
\DefineVerbatimEnvironment{Highlighting}{Verbatim}{commandchars=\\\{\}}
% Add ',fontsize=\small' for more characters per line
\newenvironment{Shaded}{}{}
\newcommand{\AlertTok}[1]{\textcolor[rgb]{1.00,0.00,0.00}{\textbf{#1}}}
\newcommand{\AnnotationTok}[1]{\textcolor[rgb]{0.38,0.63,0.69}{\textbf{\textit{#1}}}}
\newcommand{\AttributeTok}[1]{\textcolor[rgb]{0.49,0.56,0.16}{#1}}
\newcommand{\BaseNTok}[1]{\textcolor[rgb]{0.25,0.63,0.44}{#1}}
\newcommand{\BuiltInTok}[1]{\textcolor[rgb]{0.00,0.50,0.00}{#1}}
\newcommand{\CharTok}[1]{\textcolor[rgb]{0.25,0.44,0.63}{#1}}
\newcommand{\CommentTok}[1]{\textcolor[rgb]{0.38,0.63,0.69}{\textit{#1}}}
\newcommand{\CommentVarTok}[1]{\textcolor[rgb]{0.38,0.63,0.69}{\textbf{\textit{#1}}}}
\newcommand{\ConstantTok}[1]{\textcolor[rgb]{0.53,0.00,0.00}{#1}}
\newcommand{\ControlFlowTok}[1]{\textcolor[rgb]{0.00,0.44,0.13}{\textbf{#1}}}
\newcommand{\DataTypeTok}[1]{\textcolor[rgb]{0.56,0.13,0.00}{#1}}
\newcommand{\DecValTok}[1]{\textcolor[rgb]{0.25,0.63,0.44}{#1}}
\newcommand{\DocumentationTok}[1]{\textcolor[rgb]{0.73,0.13,0.13}{\textit{#1}}}
\newcommand{\ErrorTok}[1]{\textcolor[rgb]{1.00,0.00,0.00}{\textbf{#1}}}
\newcommand{\ExtensionTok}[1]{#1}
\newcommand{\FloatTok}[1]{\textcolor[rgb]{0.25,0.63,0.44}{#1}}
\newcommand{\FunctionTok}[1]{\textcolor[rgb]{0.02,0.16,0.49}{#1}}
\newcommand{\ImportTok}[1]{\textcolor[rgb]{0.00,0.50,0.00}{\textbf{#1}}}
\newcommand{\InformationTok}[1]{\textcolor[rgb]{0.38,0.63,0.69}{\textbf{\textit{#1}}}}
\newcommand{\KeywordTok}[1]{\textcolor[rgb]{0.00,0.44,0.13}{\textbf{#1}}}
\newcommand{\NormalTok}[1]{#1}
\newcommand{\OperatorTok}[1]{\textcolor[rgb]{0.40,0.40,0.40}{#1}}
\newcommand{\OtherTok}[1]{\textcolor[rgb]{0.00,0.44,0.13}{#1}}
\newcommand{\PreprocessorTok}[1]{\textcolor[rgb]{0.74,0.48,0.00}{#1}}
\newcommand{\RegionMarkerTok}[1]{#1}
\newcommand{\SpecialCharTok}[1]{\textcolor[rgb]{0.25,0.44,0.63}{#1}}
\newcommand{\SpecialStringTok}[1]{\textcolor[rgb]{0.73,0.40,0.53}{#1}}
\newcommand{\StringTok}[1]{\textcolor[rgb]{0.25,0.44,0.63}{#1}}
\newcommand{\VariableTok}[1]{\textcolor[rgb]{0.10,0.09,0.49}{#1}}
\newcommand{\VerbatimStringTok}[1]{\textcolor[rgb]{0.25,0.44,0.63}{#1}}
\newcommand{\WarningTok}[1]{\textcolor[rgb]{0.38,0.63,0.69}{\textbf{\textit{#1}}}}
\setlength{\emergencystretch}{3em} % prevent overfull lines
\providecommand{\tightlist}{%
  \setlength{\itemsep}{0pt}\setlength{\parskip}{0pt}}
\setcounter{secnumdepth}{-\maxdimen} % remove section numbering
\ifLuaTeX
\usepackage[bidi=basic]{babel}
\else
\usepackage[bidi=default]{babel}
\fi
\babelprovide[main,import]{italian}
% get rid of language-specific shorthands (see #6817):
\let\LanguageShortHands\languageshorthands
\def\languageshorthands#1{}
\ifLuaTeX
  \usepackage{selnolig}  % disable illegal ligatures
\fi
\usepackage{bookmark}
\IfFileExists{xurl.sty}{\usepackage{xurl}}{} % add URL line breaks if available
\urlstyle{same}
\hypersetup{
  pdflang={it},
  hidelinks,
  pdfcreator={LaTeX via pandoc}}

\author{}
\date{}

\begin{document}

\section{DDTA research work plan after documentation-layer
closure}\label{ddta-research-work-plan-after-documentation-layer-closure}

\textbf{WORK PLAN - REVISION 1}

Baseline: \texttt{9a04bc5b98b7f4a2b8400f2b210e7cea6b04d1a9}

This plan \textbf{supersedes earlier sequencing notes} that placed
formal Base Analysis before SecurityRequirement design. It changes the
research sequence, not the closed MR/Decision/FR semantic rules.
SecurityRequirement semantics will be defined method-neutrally before
Base Analysis/overlay construction; analytical provenance from Finding
to SEC will be integrated later.

\subsection{1. Current checkpoint}\label{current-checkpoint}

The documentation-authoring metamodel is frozen for the current thesis
scope through:

\begin{Shaded}
\begin{Highlighting}[]
\NormalTok{Project problem framing [method precondition]}
\NormalTok{        {-}\textgreater{} MacroRequirement}
\NormalTok{        {-}\textgreater{} Decision}
\NormalTok{        {-}\textgreater{} FunctionalRequirement}
\NormalTok{        {-}\textgreater{} SpecializedRequirement}
\end{Highlighting}
\end{Shaded}

Closed working rules include:

\begin{itemize}
\tightlist
\item
  regular \texttt{MR\ -\textgreater{}\ Decision\ -\textgreater{}\ FR}
  hierarchy;
\item
  exactly one MR parent per Decision;
\item
  exactly one Decision parent per FR;
\item
  controlled necessity/default Decision for rare singleton-solution
  cases;
\item
  no Decision below FR;
\item
  normative FR prose is primary and is supported by governed SPO
  references;
\item
  cross-MR capability consumption does not transfer ownership;
\item
  exact analytical representation of referenced subjects/services
  remains deferred.
\end{itemize}

The documentation baseline is reopenable only for a concrete
counterexample from the next phases.

\subsection{2. Immediate thesis-writing
checkpoint}\label{immediate-thesis-writing-checkpoint}

\subsubsection{W0 - Chapter 4 candidate}\label{w0---chapter-4-candidate}

Produce and review a standalone thesis Chapter 4 covering only:

\begin{itemize}
\tightlist
\item
  problem framing;
\item
  MR metamodel and authoring rules;
\item
  Decision metamodel, responsibility boundary and necessity/default
  rule;
\item
  FR operational semantics, normative prose + SPO, FunctionalPredicate
  and parameter classification;
\item
  cross-MR service consumption and tool-authority boundary;
\item
  validation history and explicit reopen conditions.
\end{itemize}

Do not integrate the chapter into the master thesis file until human
review is complete.

\textbf{Exit gate:} chapter accepted as an accurate description of the
current documentation baseline.

\subsection{3. Specialized Requirement and Security Requirement before
threat
methods}\label{specialized-requirement-and-security-requirement-before-threat-methods}

\subsubsection{S1 - Minimal SpecializedRequirement
semantics}\label{s1---minimal-specializedrequirement-semantics}

Define the smallest method-independent meaning shared by specialized
requirements.

Questions:

\begin{enumerate}
\def\labelenumi{\arabic{enumi}.}
\tightlist
\item
  What does specialization add to an FR: constraint, additional
  obligation, applicability condition, or a combination?
\item
  Is the relation conceptually \texttt{specializes},
  \texttt{constrains}, \texttt{appliesTo}, or something smaller/more
  neutral?
\item
  Can one FR have multiple independent specialized requirements?
\item
  How is effective obligation constructed without copying specialized
  text into the FR?
\item
  Which common identity/lifecycle/provenance concepts are inherited from
  GovernedDocument rather than redefined?
\end{enumerate}

\textbf{Exit gate:} a minimal common contract that does not assume
Security, Governance, Privacy or Performance-specific semantics.

\subsubsection{S2 - SecurityRequirement semantic
model}\label{s2---securityrequirement-semantic-model}

Define \texttt{SecurityRequirement} \textbf{before} STRIDE, STRIDE-AI or
another threat method is used as construction evidence.

Fundamental research question candidate:

\begin{quote}
Which additional security obligation must the governed functional
capability satisfy so that a required security property is preserved
under the relevant conditions, independently from the method that
discovered or motivated the obligation?
\end{quote}

Required tests:

\begin{itemize}
\tightlist
\item
  meaningful with all STRIDE category names removed;
\item
  meaningful if a different security method, expert review, incident,
  penetration test or regulatory assessment motivates it;
\item
  clearly distinguishable from FR, Decision, control/implementation
  mechanism and verification evidence;
\item
  does not embed threat-method payload in the SEC semantic core;
\item
  supports later analytical provenance without requiring that provenance
  to define what the SEC means.
\end{itemize}

\subsubsection{S3 - Independent SecurityRequirement authoring
corpus}\label{s3---independent-securityrequirement-authoring-corpus}

Author SecurityRequirement examples without running STRIDE.

Use at least two domains already understood by the study, but do not
derive the corpus from threat-category enumeration. Include positive and
negative classification examples:

\begin{itemize}
\tightlist
\item
  valid SEC;
\item
  statement that is still an FR;
\item
  statement that is a Decision;
\item
  implementation/control mechanism;
\item
  verification/evidence statement.
\end{itemize}

\textbf{Exit gate:} SEC semantics survive independent review and
regression without methodology vocabulary.

\subsection{4. Formal Base Analysis
phase}\label{formal-base-analysis-phase}

\subsubsection{BA1 - Base Analysis responsibility and
boundary}\label{ba1---base-analysis-responsibility-and-boundary}

Freeze what Base Analysis must provide and must not provide.

Candidate direction to test:

\begin{quote}
a governed, method-neutral graph of stable analysis-relevant
system/project subjects and relations, integrated from explicit governed
documentation and reviewed analytical additions.
\end{quote}

Do not use STRIDE categories to define the common graph.

\subsubsection{BA2 - BAE ontology}\label{ba2---bae-ontology}

Test and close the smallest generic subject vocabulary needed by real
project descriptions. Candidate pressure points include:

\begin{itemize}
\tightlist
\item
  Actor;
\item
  Component / Service capability;
\item
  DataResource;
\item
  DataStore;
\item
  DataFlow / Interaction;
\item
  CommunicationChannel;
\item
  InterfaceContract;
\item
  Boundary.
\end{itemize}

Open distinctions must be tested rather than copied from ThreatForge.

\subsubsection{BA3 - Base Analysis relations, provenance and
authority}\label{ba3---base-analysis-relations-provenance-and-authority}

Formalize:

\begin{itemize}
\tightlist
\item
  stable identity;
\item
  source/target/carries and other generic relations;
\item
  origin: governed document vs reviewed analytical addition;
\item
  authoritative source vs supporting sources;
\item
  lifecycle/staleness hooks;
\item
  document/SPO reference -\textgreater{} BAE resolution;
\item
  cross-MR service/capability consumption target.
\end{itemize}

\subsubsection{BA4 - Authoring-to-Base-Analysis
integration}\label{ba4---authoring-to-base-analysis-integration}

Define how MR/Decision/FR authoring creates or references
analysis-relevant system facts without automatic noun-to-BAE
canonization.

\textbf{Exit gate:} a method-neutral Base Analysis graph can be built
from governed documentation plus explicit reviewed additions without
changing the documentation hierarchy.

\subsection{5. Generic analysis envelope before concrete
overlay}\label{generic-analysis-envelope-before-concrete-overlay}

\subsubsection{A1 - AnalysisRecord common
contract}\label{a1---analysisrecord-common-contract}

Reconfirm method-independent fields for an analysis
execution/interpretation while keeping method-specific payload opaque
and plugin-owned.

\subsubsection{A2 - Common Finding
contract}\label{a2---common-finding-contract}

Reconfirm a reviewed common downstream finding envelope that can
identify affected governed subjects and preserve evidence/rationale
without copying method-specific taxonomies to the top level.

\subsubsection{A3 - Finding -\textgreater{} SecurityRequirement
integration}\label{a3---finding---securityrequirement-integration}

Only after S2/S3 and A1/A2 are stable, define analytical provenance:

\begin{Shaded}
\begin{Highlighting}[]
\NormalTok{AnalysisRecord {-}\textgreater{} Finding {-}\textgreater{} accepted/reviewed {-}\textgreater{} SecurityRequirement}
\end{Highlighting}
\end{Shaded}

The test is whether a Finding can propose/justify an SEC that conforms
to the \textbf{already defined} SEC semantics. If the SEC model must be
changed merely to fit a method payload, treat that as methodology
contamination unless a method-independent counterexample justifies the
change.

\subsection{6. Method-neutral overlay/plugin
contract}\label{method-neutral-overlayplugin-contract}

\subsubsection{O1 - Generic overlay
contract}\label{o1---generic-overlay-contract}

Define what every analysis method/plugin may receive and return without
prescribing a specific taxonomy:

\begin{itemize}
\tightlist
\item
  compatibility/version contract;
\item
  scope selection over Base Analysis subjects/relations/FRs;
\item
  method-owned payload schema;
\item
  diagnostics for missing/inadequate analysis context;
\item
  deterministic/common Finding derivation boundary;
\item
  failure/incompatibility behavior;
\item
  preservation of common records even if a plugin is unavailable.
\end{itemize}

\textbf{Exit gate:} a hypothetical non-STRIDE security method can
implement the contract without changing the common core.

\subsection{7. Concrete methodology
tests}\label{concrete-methodology-tests}

\subsubsection{O2 - STRIDE overlay}\label{o2---stride-overlay}

Implement/use STRIDE as the first concrete overlay only after S1-S3,
BA1-BA4 and O1 are frozen enough for falsification.

Evaluation questions:

\begin{itemize}
\tightlist
\item
  Does STRIDE consume Base Analysis without method-specific changes to
  BAE?
\item
  Can all STRIDE-specific semantics remain in the plugin payload/logic?
\item
  Can findings be expressed in the common Finding envelope?
\item
  Can accepted findings produce SEC instances without changing the SEC
  semantic core?
\end{itemize}

\subsubsection{O3 - Closed-loop design
improvement}\label{o3---closed-loop-design-improvement}

Exercise the complete DDTA loop:

\begin{Shaded}
\begin{Highlighting}[]
\NormalTok{governed documentation}
\NormalTok{ {-}\textgreater{} Base Analysis}
\NormalTok{ {-}\textgreater{} STRIDE overlay}
\NormalTok{ {-}\textgreater{} AnalysisRecord}
\NormalTok{ {-}\textgreater{} reviewed Finding}
\NormalTok{ {-}\textgreater{} SecurityRequirement}
\NormalTok{ {-}\textgreater{} effective obligation (FR + SEC)}
\NormalTok{ {-}\textgreater{} Decision/FR/design revision where required}
\NormalTok{ {-}\textgreater{} Base Analysis refresh}
\NormalTok{ {-}\textgreater{} re{-}analysis}
\end{Highlighting}
\end{Shaded}

The goal is not merely threat enumeration. The experiment must show
whether analysis produces a governed security obligation that changes
design/documentation and whether re-analysis can observe the changed
state.

\subsubsection{O4 - STRIDE-AI second implementation/evaluation
method}\label{o4---stride-ai-second-implementationevaluation-method}

Use STRIDE-AI as the second in-scope method for thesis
implementation/evaluation. It must reuse the same common Base Analysis,
AnalysisRecord, Finding and SecurityRequirement contracts.

\subsection{8. Method-neutrality challenges outside primary
implementation
scope}\label{method-neutrality-challenges-outside-primary-implementation-scope}

After STRIDE and STRIDE-AI, use LINDDUN and optionally one additional
method as \textbf{neutrality challenges}, not as required
implementation/evaluation methods.

Purpose:

\begin{itemize}
\tightlist
\item
  identify assumptions accidentally specific to security-threat
  taxonomies;
\item
  test whether Base Analysis subjects/relations are sufficiently
  generic;
\item
  test whether the overlay contract can describe different analytical
  semantics;
\item
  do not claim full methodology neutrality until these challenges are
  performed.
\end{itemize}

\subsection{9. ThreatForge role}\label{threatforge-role}

ThreatForge remains a reference implementation/tooling case study. The
sequence above first defines DDTA contracts, then checks which contracts
ThreatForge supports or must be changed to support.

Do not let current ThreatForge schemas establish the semantic truth of
SecurityRequirement, Base Analysis or overlay contracts.

\subsection{10. Reopen policy}\label{reopen-policy}

Reopen the \texttt{MR\ -\textgreater{}\ Decision\ -\textgreater{}\ FR}
documentation baseline only if a later phase produces a concrete
recurring semantic failure, for example:

\begin{itemize}
\tightlist
\item
  SEC cannot constrain/specialize FR without false duplication or
  ownership;
\item
  Base Analysis cannot represent facts required by the documented
  behavior without changing authoring semantics;
\item
  overlay input requires method-specific data to be embedded in
  MR/Decision/FR;
\item
  the constant Decision layer produces a real semantic contradiction
  rather than the cost of a rare necessity/default node.
\end{itemize}

Any reopening must name the counterexample and regress all previously
accepted corpora.

\subsection{11. Next authorized
microstep}\label{next-authorized-microstep}

After Chapter 4 review, start \textbf{S1: Minimal SpecializedRequirement
semantics}, then \textbf{S2: SecurityRequirement semantics}, before Base
Analysis or STRIDE construction work.

\end{document}
